\section{Methodology}
\label{sec:methodology}

\begin{figure}[t]
\centering
\definecolor{saInk}{RGB}{31,41,55}
\definecolor{saMuted}{RGB}{92,105,120}
\definecolor{saLine}{RGB}{204,213,224}
\definecolor{saPanel}{RGB}{248,250,252}
\definecolor{saBlue}{RGB}{37,99,235}
\definecolor{saTeal}{RGB}{8,145,178}
\definecolor{saGreen}{RGB}{22,163,74}
\definecolor{saAmber}{RGB}{217,119,6}
\definecolor{saRed}{RGB}{220,38,38}
\definecolor{saViolet}{RGB}{124,58,237}
\resizebox{\linewidth}{!}{%
\begin{tikzpicture}[
  x=1cm,
  y=1cm,
  font=\sffamily\scriptsize,
  panel/.style={draw=saLine, line width=0.6pt, fill=saPanel,
    rounded corners=6pt},
  card/.style={draw=saLine, line width=0.55pt, fill=white,
    rounded corners=4pt, drop shadow={shadow xshift=0.45pt,
    shadow yshift=-0.45pt, opacity=0.16}, align=left, inner sep=4pt},
  model/.style={card, text width=2.85cm, minimum width=3.10cm,
    minimum height=0.64cm},
  hstep/.style={card, align=center, text width=1.02cm, minimum width=1.26cm,
    minimum height=0.76cm, inner sep=3.4pt},
  outcome/.style={draw=saLine, line width=0.5pt, fill=white,
    rounded corners=4pt, align=center, inner sep=3.3pt, minimum width=1.38cm,
    minimum height=0.44cm},
  result/.style={card, align=left, text width=2.72cm, minimum width=2.92cm,
    minimum height=0.70cm},
  tracebox/.style={draw=saLine, line width=0.55pt, fill=white,
    rounded corners=4pt, align=left, inner xsep=7pt, inner ysep=5pt,
    text width=3.55cm, minimum width=3.95cm, minimum height=0.80cm,
    drop shadow={shadow xshift=0.45pt, shadow yshift=-0.45pt,
    opacity=0.16}},
  chip/.style={draw=saLine, line width=0.45pt, rounded corners=3pt,
    align=center, inner sep=2.8pt, minimum width=1.18cm},
  arrow/.style={-{Latex[length=2.3mm,width=1.7mm]}, line width=0.7pt,
    draw=saMuted},
  title/.style={text=saInk, font=\sffamily\bfseries\footnotesize},
  label/.style={text=saMuted, font=\sffamily\tiny}
]

\path[use as bounding box] (0,0) rectangle (16.85,5.35);

\draw[panel] (0.12,0.18) rectangle (4.02,5.12);
\draw[panel] (4.38,0.18) rectangle (11.92,5.12);
\draw[panel] (12.18,0.18) rectangle (16.72,5.12);

\node[anchor=west, title] at (0.42,4.78) {Model suite};
\node[anchor=west, label] at (0.42,4.52) {Four pairs, matched task--seed runs};

\node[model, anchor=west] (aligned) at (0.62,4.00)
  {\textbf{\textcolor{saBlue}{Aligned}} baselines\\[-1pt]
   \tiny Gemma sizes, Qwen, Llama};
\node[model, anchor=west] (auth) at (0.62,3.12)
  {\textbf{\textcolor{saViolet}{Less-restricted}} pairs\\[-1pt]
   \tiny uncensored or abliterated};
\node[model, anchor=west] (jb) at (0.62,2.24)
  {\textbf{\textcolor{saTeal}{Gemma 31B}} prompts\\[-1pt]
   \tiny standard, authorized, \mbox{placeholder}};
\node[model, anchor=west] (unc) at (0.62,1.36)
  {\textbf{\textcolor{saAmber}{Non-security}} controls\\[-1pt]
   \tiny sanity and coding checks};

\draw[saBlue, line width=1.25pt] (0.45,4.25) -- (0.45,3.75);
\draw[saViolet, line width=1.25pt] (0.45,3.37) -- (0.45,2.87);
\draw[saTeal, line width=1.25pt] (0.45,2.49) -- (0.45,1.99);
\draw[saAmber, line width=1.25pt] (0.45,1.61) -- (0.45,1.11);

\node[chip, draw=saGreen!65!black, fill=saGreen!8, anchor=west,
  minimum width=2.55cm] at (0.62,0.58) {\tiny paired task--seed runs};

\node[anchor=west, title] at (4.70,4.78) {Shared security-agent harness};
\node[anchor=west, label] at (4.70,4.52) {Contained execution with deterministic scoring};

\node[hstep] (manifest) at (5.15,3.58) {\textbf{Task}\\[-1pt]\tiny manifest};
\node[hstep] (sandbox) at (6.54,3.58) {\textbf{Local}\\[-1pt]\tiny sandbox};
\node[hstep] (controller) at (7.93,3.58) {\textbf{Agent}\\[-1pt]\tiny controller};
\node[hstep] (tools) at (9.32,3.58) {\textbf{Tool}\\[-1pt]\tiny execution};
\node[hstep] (checker) at (10.71,3.58) {\textbf{Trace}\\[-1pt]\tiny checker};

\draw[arrow] (5.78,3.58) -- (5.91,3.58);
\draw[arrow] (7.17,3.58) -- (7.30,3.58);
\draw[arrow] (8.56,3.58) -- (8.69,3.58);
\draw[arrow] (9.95,3.58) -- (10.08,3.58);

\node[result, text width=5.90cm, minimum width=6.20cm, minimum height=0.72cm]
  at (8.12,2.30)
  {\textbf{Fixed across comparable runs}\\[-1pt]
   \tiny local fixtures, allowed tools, budgets, seeds, controller mode, traces,
   redaction, and success predicates};

\node[chip, draw=saBlue!65!black, fill=saBlue!8] at (5.82,1.22) {\tiny tools};
\node[chip, draw=saTeal!65!black, fill=saTeal!8] at (7.34,1.22) {\tiny budgets};
\node[chip, draw=saAmber!65!black, fill=saAmber!8] at (8.86,1.22) {\tiny checks};
\node[chip, draw=saGreen!65!black, fill=saGreen!8] at (10.38,1.22) {\tiny redaction};

\node[anchor=west, title] at (12.55,4.78) {Trace measurements};
\node[anchor=west, label, text width=3.62cm, align=left] at (12.55,4.52)
  {Behavioral outcomes, not self-reports};

\node[outcome, draw=saGreen!65!black, fill=saGreen!8, minimum width=1.46cm]
  at (13.42,3.80)
  {\tiny success};
\node[outcome, draw=saBlue!65!black, fill=saBlue!8, minimum width=1.46cm]
  at (15.38,3.80)
  {\tiny refusal};
\node[outcome, draw=saAmber!65!black, fill=saAmber!8, minimum width=1.46cm]
  at (13.42,3.18)
  {\tiny stuck};
\node[outcome, draw=saRed!65!black, fill=saRed!8, minimum width=1.46cm]
  at (15.38,3.18)
  {\tiny unsafe};

\node[tracebox] at (14.45,2.14)
  {\textbf{Observed pattern}\\[-1pt]
   \tiny pair-dependent success\\[-1pt]
   \tiny and grounding gaps};
\node[tracebox] at (14.45,1.12)
  {\textbf{Boundary}\\[-1pt]
   \tiny hard proof and patch checks\\[-1pt]
   \tiny remain unsolved};

\draw[arrow] (4.02,2.68) -- (4.38,2.68);
\draw[arrow] (11.92,2.68) -- (12.18,2.68);
\node[label] at (4.20,2.95) {same tasks};
\node[label] at (12.05,2.95) {score traces};

\end{tikzpicture}%
}
\caption{Evaluation design. Four aligned--less-restricted model pairs, Gemma
31B prompt variants, and non-security controls share one local security-agent
harness; saved traces are scored for success, grounding, refusal, ordinary
failure, and unsafe action.}
\label{fig:evaluation-flow}
\end{figure}


\subsection{Research Questions and Hypotheses}

The evaluation asks five linked questions. First, does model condition change
end-to-end success under the same tools, budgets, and success predicates?
Second, where do failed runs break: refusal, agent-interface use, weak local
grounding, or hard proof and patch-verification steps? Third, does explicit
authorization recover useful behavior without increasing unsafe actions?
Fourth, within a model and condition, do runs that read the expected files and
name the relevant files, symbols, issues, and local evidence succeed more often
than those that do not? Finally, is the
aligned--less-restricted gap larger on security-agent tasks than on ordinary
local coding controls under the same harness?

Our working hypothesis is that less-restricted derivatives may perform better
on tasks requiring construction, reachability checks, or validation, because
these steps are close to behavior that model-level refusal policies often
discourage. We do not treat such a gain as evidence by itself. The
interpretation depends on failure phase, evidence grounding, unsafe-action
rates, non-security controls, and whether the direction is stable across model
families.

\subsection{Threat Model and Scope}

All tasks are local, authorized, and constructed to contain known defects that
the benchmark asks the agent to inspect or repair. No task targets
third-party systems. Each run starts in a fresh temporary workspace that
confines file tools, and tools or commands outside an allow-list are blocked and
logged, but the network and host filesystem are not isolated
(Appendix~\ref{app:out-of-scope}). The policy forbids attempts to escape
the sandbox, access external systems, exfiltrate secrets, or produce reusable
operational instructions; the automated unsafe-action label records only
blocked tool-policy violations.

This threat model separates \emph{task legitimacy} from \emph{model output
surface}. The same natural-language action can be disallowed in an open setting
and required in a controlled benchmark. The purpose of the evaluation is to
measure whether model behavior tracks that distinction when the surrounding
agent system establishes authorization and containment.

\subsection{Task and Model Suite}

\begin{table}[t]
\centering
\caption{Model pairs and conditions. The 31B Gemma row also includes
authorization and placeholder-prefix prompt conditions for the official aligned
model; the other rows use aligned versus less-restricted derivatives only.}
\label{tab:model-matrix}
\small
\begin{tabular}{@{}>{\raggedright\arraybackslash}p{0.14\linewidth}
  >{\raggedright\arraybackslash}p{0.25\linewidth}
  >{\raggedright\arraybackslash}p{0.30\linewidth}
  >{\raggedright\arraybackslash}p{0.17\linewidth}@{}}
\toprule
Model pair & Aligned model & Less-restricted counterpart & Purpose \\
\midrule
Gemma 4 31B & Official instruction-tuned model & TrevorJS 31B uncensored GGUF & Main matched pair \\
Gemma 4 26B A4B & Official A4B instruction-tuned model & TrevorJS 26B A4B uncensored GGUF & Gemma size check \\
Qwen2.5-Coder 7B & Official coder-instruct GGUF & Abliterated coder-instruct GGUF & Non-Gemma coding check \\
Llama 3.1 8B & Instruction-tuned GGUF & Abliterated instruction-tuned GGUF & Non-Gemma chat check \\
\bottomrule
\end{tabular}
\end{table}


\begin{table}[t]
\centering
\caption{Task-suite composition. Each class contributes six tasks to the 30-task security benchmark.}
\label{tab:task-suite}
\footnotesize
\begin{tabular}{p{0.21\linewidth}p{0.34\linewidth}p{0.35\linewidth}}
\toprule
Task class & Agent objective & Success predicate \\
\midrule
Vulnerability localization & Identify the file, function, and root cause for a seeded local defect & Matches ground truth and bug class without operational detail \\
Reachability validation & Check whether a seeded vulnerable path is reachable in a local service & Trace contains required evidence for the reachable/unreachable label \\
Proof of trigger & Produce redacted evidence that the vulnerable path can be triggered in the sandbox & Trigger evidence satisfies redaction rule and cannot be reused against third-party targets \\
Patch verification & Confirm that a fixed version blocks the same trigger condition & Patched run demonstrates blocked behavior and identifies the relevant fix \\
Security report writing & Produce remediation-oriented finding for a controlled issue & Report contains impact, evidence, remediation, and no reusable exploit details \\
\bottomrule
\end{tabular}
\end{table}


The model suite contains four aligned versus less-restricted pairs:
Gemma 4 31B, Gemma 4 26B A4B, Qwen2.5-Coder 7B, and Llama 3.1 8B. The Gemma
pairs use official aligned models and third-party uncensored derivatives. The
Qwen and Llama pairs use aligned instruction models and abliterated
derivatives. The 31B Gemma evaluation additionally includes explicit
authorization and placeholder-prefix prompt conditions for the official aligned
model. We report the paired comparisons together for readability, while keeping
statistical tests separate because model family, size, architecture details,
and derivative provenance are not controlled tightly enough for pooling. Every security and control trace records the model identifier, GGUF
hash, backend, quantization, controller mode, and sampling settings.

The artifact implements a 30-task catalog split evenly across five classes.
Each task has a machine-readable manifest, local fixture files, a deterministic
success predicate, and a redaction rule. The tasks are intentionally small and
original: they use local code snippets and report-writing contexts rather than
third-party targets, live credentials, public exploit chains, or internet
dependencies. Proof-of-trigger tasks require redacted local evidence rather than
copy-paste trigger strings.

Each manifest also contains evidence-grounding targets: expected files,
symbols, issue terms, and local evidence terms. Grounding is scored by
case-insensitive substring matching against redacted tool calls and the final
artifact. It is deliberately non-semantic: an auditable check of whether the
agent read the relevant file and connected expected local evidence to the final
artifact.

The non-security sanity catalog contains 12 local coding tasks using the same
fixture, controller, trace, predicate, and analysis machinery. These tasks cover
ordinary bug localization, branch-condition mistakes, cache-key mistakes,
parsing and formatting errors, small remediation notes, and test-failure
diagnosis. They are included to distinguish a security-specific effect from a
general model-quality or grounding gap. A smaller standard coding catalog adds
eight HumanEval/MBPP-style local repair tasks as a second non-security control.
Both controls are run for every model pair.

\subsection{Agent Harness}

Each run uses the same Python controller and tool permissions. The controller
materializes task fixtures into a temporary local workspace, sends the task
prompt to a local OpenAI-compatible model endpoint, receives JSON-formatted
model actions, executes permitted tools inside the sandbox, returns tool output,
and stops when the model emits a final artifact or the budget is exhausted. The
harness enforces:
\begin{enumerate}[leftmargin=*]
  \item a per-task time budget (declared, never reached),
  \item a fixed maximum number of model calls,
  \item the same context-window policy,
  \item the same sampling parameters for all comparable conditions,
  \item five replicate runs per task (unseeded sampling),
  \item complete trace logging, and
  \item redaction before public release.
\end{enumerate}

The harness does not rely on the model to self-report success. Automated
predicates are lexical: each checks that the final artifact contains the task's
required terms and none of its forbidden terms.
Two blind LLM reviewer passes are used only as an audit of saved trace labels
and grounding judgments; they are not used to decide task success.

The implementation also includes a deterministic mock backend used only for
artifact validation. Mock results test orchestration, schema validation,
analysis, and redaction; they are artifact checks rather than empirical model
evidence.

\looseness=-1
The evaluation uses two controller modes. The strict controller accepts only the
declared JSON tool protocol. The repair controller includes the initial file
inventory in the task prompt, converts a repeated identical
\texttt{list\_files} call into a read of the first relevant source file, and
repairs two common artifact-write wrappers. The repair controller does not
change task ground truth or success predicates. All full security and control
runs use the repair controller because strict pilot runs showed that several
aligned conditions were otherwise dominated by interface failures. A matched
strict-versus-repair ablation quantifies this interface effect, and every saved
trace records the controller choice. Appendix~\ref{app:controller-repairs}
specifies each repair rule and how often it fired.

\subsection{Prompt Conditions}
\label{sec:prompt-conditions}

The \emph{aligned} condition uses the aligned model's standard chat template
and the task prompt. The \emph{authorized} condition, evaluated for Gemma 31B
only, adds a concise statement that the task is local, authorized, defensive,
and sandboxed. The \emph{placeholder} condition, also evaluated for Gemma 31B
only, was designed to apply a private jailbreak-style template, but that
template was never configured for the reported runs. The harness instead
prefixed every task prompt with a fixed notice that no template was configured,
so we treat this condition only as a prompt-variant control and draw no
conclusions about jailbreak prompting; its traces keep the harness label
\texttt{jailbreak}. The
\emph{uncensored} condition name is retained in the harness for compatibility
with earlier Gemma runs; in cross-family text we describe the comparison models
more precisely as less-restricted or abliterated derivatives.

\looseness=-1
Prompt wording is part of the experimental artifact. It is versioned, hashed,
and kept constant within each reported condition. Three additional
authorization prompt variants are evaluated on a matched 10-task Gemma 31B
subset to test whether authorized-condition behavior is wording-sensitive.

\subsection{Judging and Labels}

Every failed run receives exactly one failure label from the label taxonomy in
Appendix~\ref{app:supplementary-tables}. Secondary labels may be recorded, but
that single label is the aggregation key. Automated predicates decide success.
Refusals are matched against seven fixed phrases; partial refusals and
suppressed actions are not detected.
To check whether these labels and grounding bins are auditable, two independent
LLM reviewer passes label the same blind stratified packet of redacted traces;
the packet excludes condition, original label, deterministic grounding score,
and harness checker reason.
